\newpage
\section{等差数列与等比数列}
\label{sec:analysis-arithmetic-geometric}

\subsection{等差数列}

\begin{definition}
若数列相邻两项的差\(a_{n+1}-a_n=d\)与\(n\)无关，称为{\heiti{}等差数列}，常数\(d\)称为公差。
\end{definition}

将\(a_2-a_1=d,\ldots,a_n-a_{n-1}=d\)相加，中间各项消去，得到
\begin{equation}\label{eq:analysis-arithmetic-term}
 a_n=a_1+(n-1)d.
\end{equation}
任意两项满足\(a_n-a_m=(n-m)d\)；下标和相同的两对项，其和也相同。将前\(n\)项正序、逆序各写一次，逐项相加，每一对都是\(a_1+a_n\)，所以
\begin{equation}\label{eq:analysis-arithmetic-sum}
 S_n=\frac{n(a_1+a_n)}2=na_1+\frac{n(n-1)}2d.
\end{equation}
特别地，\(1+2+\cdots+n=n(n+1)/2\)。等差数列的部分和是关于\(n\)的不超过二次的多项式，且常数项为零；反过来，由\(a_n=S_n-S_{n-1}\)也可以恢复等差数列。

\begin{example}
等差数列满足\(a_3+a_7=20,S_{10}=110\)。求通项与前\(n\)项和，再求使\(S_n>1000\)成立的最小正整数\(n\)。
\end{example}

\subsection{等比数列}

\begin{definition}
一个各项非零的数列，若相邻两项之比\(a_{n+1}/a_n=q\)与\(n\)无关，称为{\heiti{}等比数列}，常数\(q\)称为公比。
\end{definition}

将相邻两项的比连乘，得到
\begin{equation}\label{eq:analysis-geometric-term}
 a_n=a_1q^{n-1}.
\end{equation}
它描述固定倍率的变化；\(q<0\)时符号交替，因而固定倍率不一定产生单调数列。

设\(S_n=a_1+a_1q+\cdots+a_1q^{n-1}\)，乘以\(q\)后各项依次错开一位。两式相减，中间项消去，得\((1-q)S_n=a_1(1-q^n)\)，因而
\begin{equation}\label{eq:analysis-geometric-sum}
 S_n=a_1\frac{1-q^n}{1-q}\quad(q\ne1);
 \qquad S_n=na_1\quad(q=1).
\end{equation}
这种乘以公比使两组项对齐，再消去共同部分的方法，称为{\heiti{}错位相减}。

\begin{example}
等比数列满足\(a_2=6,a_5=162\)，求首项、公比及前\(n\)项和。另取\(b_n=(-1/2)^{n-1}\)，写出其部分和，比较原数列与部分和数列的变化。
\end{example}
